problem:  
Let \( (M,g) \) be a compact Riemann surface and \( (N,h) \) a compact Riemannian manifold. Consider the harmonic map flow  
\[
\partial_t u = \tau(u), \qquad u: M\times [0,T)\to N,
\]
with smooth initial data \(u_0\).  
Assume the flow develops a finite-time singularity at \(T<\infty\).  Denote by \(\lambda(t) = \bigl(\!\sup_{x\in M} |\nabla u(x,t)|\bigr)^{-1}\)  the *concentration scale*.  
Suppose after rescaling near the singular point \(x_0\in M\),
\[
u_k(x) := u\big(x_0 + \lambda(t_k)x, t_k\big),
\]
for some sequence \(t_k\uparrow T\), one obtains a limiting harmonic sphere \(\phi: S^2\to N\).

Let \(L_\phi\) denote the Jacobi operator associated with \(\phi\) on \(\phi^*TN\), and let \(\sigma_-(L_\phi)\subset(-\infty,0)\) be its negative spectrum.  

**Question:**  
For generic metrics \(h\) on \(N\), to what extent are the asymptotic dynamics of the blow-up scale \(\lambda(t)\) determined by the negative spectrum of \(L_\phi\)?  In particular:

1. *(Spectral–dynamic correspondence)* Does there exist a functional relation—possibly implicit—between the unstable Jacobi eigenvalues \(\sigma_-(L_\phi)\) and the asymptotic evolution law of \(\lambda(t)\) near \(t = T\)?  That is, can the leading-order modulation of \(\lambda(t)\) be expressed in terms of the most unstable eigenmode(s) of \(L_\phi\)?  
2. *(Spectral rigidity of singular profiles)* For two harmonic spheres \(\phi_1,\phi_2:S^2\to N\) with distinct negative spectral data (e.g., different Morse indices or non‐isomorphic Jacobi spectra), is it possible for the same harmonic map flow to produce both as bubbles appearing at comparable spatio–temporal scales?  
3. *(Spectral stratification)* Do finite-time singularities of the harmonic map flow admit a discrete classification by the negative spectra of Jacobi operators of the resulting harmonic spheres, analogous to the organization of dynamical systems by Lyapunov spectra?

In other words, is there a universal *spectral law of blow-up*, where bubble formation rates and hierarchies are discretized according to the unstable eigenmodes of their limiting harmonic spheres?

Why is it a "good" problem:  
This problem is precise in established analytic terms: it employs standard concentration scales in bubbling analysis and the well-defined Jacobi spectrum of the limiting harmonic sphere. Yet it proposes a fundamentally new conceptual bridge—linking the *spectral theory of the linearized energy functional* with the *dynamical scaling regimes* of nonlinear geometric flow singularities. Demonstrating any such spectral control would constitute a spectral–dynamical correspondence principle for harmonic map blow-up, paralleling modulation theories in soliton dynamics but in a geometric flow setting. The problem is simple to state—do the Jacobi eigenvalues dictate the rates or types of bubble formation?—but deeply nontrivial, calling for new insights into coupling between linearized instability, nonlinear modulation, and geometric rescaling. By uniting spectral geometry, nonlinear PDE, and singularity formation theory, it exhibits the clarity, depth, and cross-disciplinary potential that define a “good” research-level mathematical problem.